\chapter{Gaussian posteriors for adaptively collected data and the unit-ball lower bound}
\label{chap:pre_bayes}

This chapter replaces the citation of \cite{chen2024assouad} in Appendix~H of the paper by a
self-contained argument: for every (adaptive, possibly randomized) identification algorithm
with budget $T$ on the unit ball $\ball_d$, there is a reward vector $\theta$ of norm
$O(d/\sqrt T)$ under which the expected simple regret is $\Omega(d/\sqrt T)$; consequently an
$(\eps,\delta)$-PAC algorithm on $\ball_d$ needs $T = \Omega(d^2/\eps^2)$ samples
(Corollary~\ref{cor:ball_pac_lower}, used in Chapter~\ref{chap:log_gains}). The proof is
Bayesian: under a Gaussian prior $\theta\sim\Normal(0,\sigma^2I_d)$, the posterior mean of
$\theta$ given the history is $m_T = (\sigma^{-2}I+S_T)^{-1}b_T$, and the recommendation cannot
correlate with $\theta$ better than with $m_T$. We never use abstract conditional expectations:
the key identity $\E\ip{\varphi}{\theta-m_T} = 0$ is proved by Fubini, the likelihood ratio of
the history with respect to the noise-only environment, and an explicit completion of the
square in $\R^d$. The likelihood ratio is first established for general stationary environments
whose kernels have densities, as the environment analogue of LML's \texttt{Algorithm.density}.

\textbf{What Mathlib/LML provides.} Measures with densities \texttt{Measure.withDensity},
\texttt{withDensity\_mul}, \texttt{Measure.compProd\_withDensity}
(\texttt{Mathlib/Probability/Kernel/CompProdEqIff.lean}: $\mu\otimes(\kappa.\mathrm{withDensity}\,f) = (\mu\otimes\kappa).\mathrm{withDensity}(f)$),
and in LML \texttt{ForMathlib/Probability/WithDensity.lean}: \texttt{compProd\_withDensity\_left},
\texttt{compProd\_withDensity\_withDensity}, \texttt{map\_equiv\_withDensity},
\texttt{Kernel.compProd\_withDensity\_left} (a density on the \emph{left} factor of a kernel
composition-product). The algorithm density
\texttt{Algorithm.density} and \texttt{IsAlgEnvSeq.hasLaw\_history\_withDensity}
(\texttt{SequentialLearning/AlgorithmDensity.lean}) are the model for
Lemma~\ref{lem:env_likelihood_ratio}. Bayesian sequences: \texttt{IsBayesAlgEnvSeq},
\texttt{bayesStationaryEnv}, \texttt{IT.bayesTrajMeasure}, \texttt{IsBayesAlgEnvSeq.ae\_IsAlgEnvSeq}
(\texttt{SequentialLearning/BayesStationaryEnv.lean}) give the joint law of a parameter drawn
from a prior and the trajectory. Gaussians: \texttt{gaussianReal\_of\_var\_ne\_zero},
\texttt{gaussianPDFReal}, \texttt{integral\_gaussian}, \texttt{multivariateGaussian},
\texttt{stdGaussian}, the linear change of variables for the Lebesgue measure
\texttt{MeasureTheory.Measure.map\_linearMap\_addHaar\_eq\_smul\_addHaar} and
\texttt{MeasureTheory.integral\_map\_equiv}.

\textbf{Formalization note.} The chapter is formalized along the plan above, with two
deviations. The posterior identities of Section~\ref{sec:pre_bayes_posterior} are obtained
from Stein's identity (Gaussian integration by parts, Lemma~\ref{lem:gauss_ibp}, extended to
functions of exponential growth in Lemma~\ref{lem:gauss_ibp_exp}) applied to the tilt
$\theta \mapsto \exp(\ip{b}{\theta} - \tfrac12\theta^\top S\theta)$, instead of the Lebesgue
density of the prior and a change of variables; no density with respect to the Lebesgue
measure on $\R^d$ is ever used. And the lower bound $\E\norm\theta \ge \sigma d/\sqrt{d+2}$ of
the outline is replaced by the sharper $\E\norm\theta \ge \sqrt{2/\pi}\,\sigma\sqrt d$ of
Lemma~\ref{lem:gauss_norm_lower}, which improves the constants (we keep the outline's
headline constants).

\section{Likelihood ratio of the history with respect to a reference environment}
\label{sec:pre_bayes_lr}

\begin{definition}[Environment density along a history]\label{def:pb_env_density}
\lean{Learning.envDensity, Learning.measurable_envDensity, Learning.envDensity_zero, Learning.envDensity_snoc}\leanok
Let $\mathcal{A},\mathcal{Y}$ be measurable spaces and $\rho : \mathcal{A}\times\mathcal{Y}\to[0,\infty]$
be measurable. For $n\ge0$ and a history $h = (x_t,y_t)_{t\le n}$ define
$D_n^\rho(h) := \prod_{t=0}^n\rho(x_t,y_t)$, a measurable function of $h$. For a history of
length $T\ge1$ we write $D_T^\rho := D_{T-1}^\rho$.
\end{definition}

\begin{lemma}[Composition-product of a kernel with a kernel having a density]
\label{lem:pb_kernel_compProd_withDensity_right}
\lean{ProbabilityTheory.Kernel.compProd_withDensity_right, ProbabilityTheory.Kernel.prodMkLeft_withDensity}\leanok
Let $\pi$ be an s-finite kernel from $\mathcal{H}$ to $\mathcal{A}$, $\eta$ an s-finite kernel
from $\mathcal{H}\times\mathcal{A}$ to $\mathcal{Y}$, and
$f : (\mathcal{H}\times\mathcal{A})\times\mathcal{Y}\to[0,\infty]$ measurable such that the
kernel $\eta.\mathrm{withDensity}(f)$ is s-finite (in all applications below it is a Markov
kernel). Then, as kernels from $\mathcal{H}$ to $\mathcal{A}\times\mathcal{Y}$,
\[
  \pi\otimes\big(\eta.\mathrm{withDensity}(f)\big)
  = (\pi\otimes\eta).\mathrm{withDensity}\big((h,(a,y))\mapsto f((h,a),y)\big).
\]
In particular, if $f((h,a),y) = \rho(a,y)$ does not depend on $h$, the density of the
composition-product is $(h,(a,y))\mapsto\rho(a,y)$, i.e.\ $\rho\circ\mathrm{snd}$.
\end{lemma}

\begin{proof}
\leanok
(The Lean proof is not fiberwise: both sides are evaluated on a measurable set through
\texttt{Kernel.compProd\_apply}, \texttt{Kernel.withDensity\_apply'} and
\texttt{Kernel.lintegral\_compProd}, exactly as Mathlib's measure-level
\texttt{Measure.compProd\_withDensity}.) Fiberwise. Fix $h\in\mathcal{H}$. For every kernel $\eta'$ from $\mathcal{H}\times\mathcal{A}$
to $\mathcal{Y}$, $(\pi\otimes\eta')(h) = \pi(h)\otimes\eta'_h$, where
$\eta'_h := \eta'((h,\cdot))$ is the section kernel from $\mathcal{A}$ to $\mathcal{Y}$
(Mathlib \texttt{Kernel.compProd\_apply\_eq\_compProd\_sectR}). The section of a kernel with
a density is the section with the sectioned density:
$(\eta.\mathrm{withDensity}(f))_h = \eta_h.\mathrm{withDensity}\big((a,y)\mapsto f((h,a),y)\big)$,
both sides being the measure $\eta((h,a)).\mathrm{withDensity}(f((h,a),\cdot))$ at each $a$
(\texttt{Kernel.withDensity\_apply}). The measure-level identity
\texttt{Measure.compProd\_withDensity} (\texttt{Mathlib/Probability/Kernel/CompProdEqIff.lean})
then gives
\[
  \pi(h)\otimes\big(\eta_h.\mathrm{withDensity}(f((h,\cdot),\cdot))\big)
  = \big(\pi(h)\otimes\eta_h\big).\mathrm{withDensity}\big((a,y)\mapsto f((h,a),y)\big),
\]
and the right-hand side is $\big((\pi\otimes\eta).\mathrm{withDensity}(\tilde f)\big)(h)$
with $\tilde f(h,(a,y)) := f((h,a),y)$, by \texttt{Kernel.withDensity\_apply} again
($\tilde f$ is measurable, being $f$ composed with a measurable rearrangement of the
variables) and \texttt{Kernel.compProd\_apply\_eq\_compProd\_sectR}. Two kernels that agree
at every point are equal (\texttt{Kernel.ext}).
\end{proof}

\begin{lemma}[Environment likelihood ratio]\label{lem:env_likelihood_ratio}
\lean{Learning.IsAlgEnvSeq.map_history_eq_withDensity_env, Learning.IsAlgEnvSeq.map_finHistory_eq_withDensity_env, Learning.toFinHistory_IicSuccProd_symm}\leanok
\uses{def:pb_env_density, lem:pb_kernel_compProd_withDensity_right}
Let $\kappa_0,\kappa$ be Markov kernels from $\mathcal{A}$ to $\mathcal{Y}$ and
$\rho : \mathcal{A}\times\mathcal{Y}\to[0,\infty]$ measurable with
$\kappa(a) = \kappa_0(a).\mathrm{withDensity}(\rho(a,\cdot))$ for every $a\in\mathcal{A}$. Let
$\mathrm{alg}$ be any algorithm and, for $n\ge0$, let $\Prob^{(n)}$ and $\Prob_0^{(n)}$ be the
laws of the history $(x_t,y_t)_{t\le n}$ under $\mathrm{stationaryEnv}(\kappa)$, resp.\
$\mathrm{stationaryEnv}(\kappa_0)$ (Definition~\ref{def:env_pair}). Then
\[
  \Prob^{(n)} = \Prob_0^{(n)}.\mathrm{withDensity}\big(D_n^\rho\big),
\]
so $\Prob^{(n)}\ll\Prob_0^{(n)}$ and, for every measurable $f\ge0$ (or $f$ integrable under
$\Prob^{(n)}$), $\int f\,d\Prob^{(n)} = \int f\,D_n^\rho\,d\Prob_0^{(n)}$.
\end{lemma}

\begin{proof}
\leanok
\uses{def:pb_env_density, lem:pb_kernel_compProd_withDensity_right}
Induction on $n$, exactly as LML's \texttt{hasLaw\_history\_withDensity} for the algorithm
density. Write $\tilde\kappa$, $\tilde\kappa_0$ for the kernels $(h,a)\mapsto\kappa(a)$,
$(h,a)\mapsto\kappa_0(a)$ (the feedback kernels of the stationary environments, LML
\texttt{feedback\_stationaryEnv}), so that
$\tilde\kappa = \tilde\kappa_0.\mathrm{withDensity}((h,a),y)\mapsto\rho(a,y))$.

\emph{Base case.} $(x_0,y_0)$ has law $p_0\otimes\kappa = p_0\otimes(\kappa_0.\mathrm{withDensity}\,\rho)
= (p_0\otimes\kappa_0).\mathrm{withDensity}((a,y)\mapsto\rho(a,y))$ (Mathlib
\texttt{Measure.compProd\_withDensity}), and $\mathrm{hist}_0$ is its image by the measurable
equivalence $\mathcal{A}\times\mathcal{Y}\simeq\mathcal{H}_0$, which commutes with
\texttt{withDensity} (LML \texttt{map\_equiv\_withDensity}); $D_0^\rho(h) = \rho(x_0,y_0)$.

\emph{Induction step.} By LML \texttt{history\_succ} and \texttt{hasCondDistrib\_step},
$\Prob^{(n+1)}$ is the image under $e^{-1} : \mathcal{H}_n\times(\mathcal{A}\times\mathcal{Y})\simeq\mathcal{H}_{n+1}$
of $\Prob^{(n)}\otimes(\pi_n\otimes\tilde\kappa)$, where $\pi_n$ is the policy kernel. Now
$\pi_n\otimes\tilde\kappa = (\pi_n\otimes\tilde\kappa_0).\mathrm{withDensity}((h,(a,y))\mapsto\rho(a,y))$
(Lemma~\ref{lem:pb_kernel_compProd_withDensity_right} with $f((h,a),y) := \rho(a,y)$; here
$\tilde\kappa = \tilde\kappa_0.\mathrm{withDensity}(f)$ is a Markov kernel), and by the
induction hypothesis
$\Prob^{(n)} = \Prob_0^{(n)}.\mathrm{withDensity}(D_n^\rho)$. A composition-product of a measure
with a density and a kernel with a density is the composition-product with the product
density (LML \texttt{compProd\_withDensity\_withDensity}):
\[
  \Prob^{(n)}\otimes(\pi_n\otimes\tilde\kappa)
  = \big(\Prob_0^{(n)}\otimes(\pi_n\otimes\tilde\kappa_0)\big).\mathrm{withDensity}\big((h,(a,y))\mapsto D_n^\rho(h)\rho(a,y)\big).
\]
Pushing forward by $e^{-1}$ (\texttt{map\_equiv\_withDensity}) gives
$\Prob^{(n+1)} = \Prob_0^{(n+1)}.\mathrm{withDensity}(D_{n+1}^\rho)$ since
$D_{n+1}^\rho(e^{-1}(h,(a,y))) = D_n^\rho(h)\rho(a,y)$. The integral formulas are
\texttt{lintegral\_withDensity\_eq\_lintegral\_mul} and
\texttt{integral\_withDensity\_eq\_integral\_smul}.
\end{proof}

\textbf{Lean remark.} \texttt{Learning.envDensity ρ n : (Fin n → 𝓐 × 𝓨) → ℝ≥0∞} is the
product $\prod_{t<n}\rho(x_t,y_t)$ over \texttt{Fin n} (so that it applies to the histories
\texttt{finHistory} used throughout), with the recursion \texttt{envDensity\_snoc}; the law
identity is proved first for LML's \texttt{history X Y n : Iic n → 𝓐 × 𝓨} by the induction of
\texttt{hasLaw\_history\_withDensity} (\texttt{map\_history\_eq\_withDensity\_env}) and
transported to \texttt{finHistory} through \texttt{toFinHistory}
(\texttt{map\_finHistory\_eq\_withDensity\_env}). The kernel-level identity
$\pi\otimes(\tilde\kappa_0.\mathrm{withDensity}\,\rho) = (\pi\otimes\tilde\kappa_0).\mathrm{withDensity}(\rho\circ\mathrm{snd})$
is Lemma~\ref{lem:pb_kernel_compProd_withDensity_right}, the right-factor analogue of LML's
\texttt{Kernel.compProd\_withDensity\_left}; it is not in Mathlib or LML and is proved
fiberwise from \texttt{Measure.compProd\_withDensity}.

\begin{lemma}[Gaussian density ratio]\label{lem:pb_gaussian_density_ratio}
\lean{ProbabilityTheory.gaussianReal_eq_withDensity_exp, ProbabilityTheory.gaussianPDFReal_eq_mul_exp, ProbabilityTheory.measurable_gaussianDensityRatio}\leanok
For $m\in\R$, $\Normal(m,1) = \Normal(0,1).\mathrm{withDensity}\big(y\mapsto e^{ym-m^2/2}\big)$.
\end{lemma}

\begin{proof}
\leanok
Both are Lebesgue measure with densities (Mathlib \texttt{gaussianReal\_of\_var\_ne\_zero}):
$\Normal(m,1) = \mathrm{Leb}.\mathrm{withDensity}(\varphi(\cdot-m))$ with
$\varphi(y) = (2\pi)^{-1/2}e^{-y^2/2}$, and
$\varphi(y-m) = \varphi(y)\,e^{ym-m^2/2}$ since $-(y-m)^2/2 = -y^2/2+ym-m^2/2$. Conclude with
\texttt{withDensity\_mul}: $(\mathrm{Leb}.\mathrm{withDensity}\,\varphi).\mathrm{withDensity}(e^{ym-m^2/2}) = \mathrm{Leb}.\mathrm{withDensity}(\varphi\cdot e^{ym-m^2/2})$.
\end{proof}

\begin{corollary}[Likelihood ratio of a linear Gaussian history]\label{cor:pb_gaussian_likelihood}
\lean{Learning.LinearBandit.stepLogLR, Learning.LinearBandit.stepLR, Learning.LinearBandit.measurable_uncurry_stepLR, Learning.LinearBandit.linearGaussianKernel_eq_withDensity, Learning.LinearBandit.likelihood, Learning.LinearBandit.likelihood_pos, Learning.LinearBandit.envDensity_stepLR, Learning.LinearBandit.measurable_likelihood, Learning.LinearBandit.measurable_likelihood_prod, Learning.IsAlgEnvSeq.map_finHistory_eq_withDensity_likelihood, Maiti2026Power.histB, Maiti2026Power.histS, Maiti2026Power.likelihood_eq_gaussTilt}\leanok
\uses{def:env, lem:env_likelihood_ratio, lem:pb_gaussian_density_ratio}
Let $\Xs$ be an action set, $\theta\in\R^d$, $T\ge1$, and for a history $h = (x_t,y_t)_{t<T}$
define
\[
  S_T(h) := \sum_{t<T}x_tx_t^\top\in\PSD,\qquad b_T(h) := \sum_{t<T}y_tx_t\in\R^d,\qquad
  \ell_\theta(h) := \exp\Big(\ip{b_T(h)}{\theta}-\tfrac12\theta^\top S_T(h)\theta\Big) .
\]
For every algorithm with actions in $\Xs$, the law $\Prob_\theta^T$ of the history of length
$T$ under $\nu_\theta$ (Definition~\ref{def:env}) satisfies
$\Prob_\theta^T = \Prob_0^T.\mathrm{withDensity}(\ell_\theta)$, where $\Prob_0^T$ is the law
under $\nu_0$ (pure noise: $y_t\sim\Normal(0,1)$). Moreover $(\theta,h)\mapsto\ell_\theta(h)$ is
jointly measurable and positive. (In Lean, \texttt{withDensity} takes an
\texttt{ℝ≥0∞}-valued density, so the density is \texttt{fun h ↦ ENNReal.ofReal (ℓ θ h)}, and
likewise $\mathrm{ENNReal.ofReal}\circ\rho_\theta$ in the proof; joint measurability of
$(\theta,h)\mapsto\mathrm{ENNReal.ofReal}(\ell_\theta(h))$ follows from that of $\ell$ and
\texttt{ENNReal.measurable\_ofReal}.) In Lean, \texttt{likelihood θ h} is the sum form
$\exp(\sum_{t<T}(y_t\ip{x_t}{\theta} - \tfrac12\ip{x_t}{\theta}^2))$, valid on any action
set; the identification with $\ip{b_T}{\theta} - \tfrac12\theta^\top S_T\theta$ is
\texttt{Maiti2026Power.likelihood\_eq\_gaussTilt} (stated on the unit ball, with
\texttt{histB}, \texttt{histS} for $b_T$, $S_T$).
\end{corollary}

\begin{proof}
\leanok
\uses{lem:env_likelihood_ratio, lem:pb_gaussian_density_ratio}
Lemma~\ref{lem:pb_gaussian_density_ratio} with $m = \ip{x}{\theta}$ gives
$\Normal(\ip{x}{\theta},1) = \Normal(0,1).\mathrm{withDensity}(\rho_\theta(x,\cdot))$ with
$\rho_\theta(x,y) := \exp(y\ip{x}{\theta}-\tfrac12\ip{x}{\theta}^2)$, jointly continuous in
$(\theta,x,y)$. Lemma~\ref{lem:env_likelihood_ratio} then gives the density
$D_T^{\rho_\theta}(h) = \prod_{t<T}\exp(y_t\ip{x_t}{\theta}-\tfrac12\ip{x_t}{\theta}^2) = \ell_\theta(h)$,
since $\sum_ty_t\ip{x_t}{\theta} = \ip{b_T}{\theta}$ and
$\sum_t\ip{x_t}{\theta}^2 = \theta^\top S_T\theta$. Since $\ell_\theta > 0$ is real-valued, the
\texttt{ℝ≥0∞} density $\mathrm{ENNReal.ofReal}(\ell_\theta)$ is finite everywhere and
$\mathrm{ENNReal.ofReal}(\prod_t\rho_\theta(x_t,y_t)) = \prod_t\mathrm{ENNReal.ofReal}(\rho_\theta(x_t,y_t))$
(\texttt{ENNReal.ofReal\_prod\_of\_nonneg}).
\end{proof}

\begin{lemma}[The law of the run as a Markov kernel in $\theta$]\label{lem:pb_kernel_of_likelihood}
\lean{Learning.LinearBandit.noiseHistLaw, Learning.LinearBandit.histKernel, Learning.LinearBandit.histKernel_apply, Learning.IsAlgEnvSeq.map_finHistory_eq_histKernel, Learning.IsAlgEnvSeq.hasLaw_finHistory_histKernel, Learning.LinearBandit.pairKernel, Learning.LinearBandit.pairKernel_apply, Learning.IdentAlg.IsRun.hasLaw_finHistory_out_pairKernel, Learning.LinearBandit.IsPAC.le_measureReal_pairKernel, Learning.LinearBandit.IsPAC.integral_simpleRegret_pairKernel_le, Maiti2026Power.expRegret, Maiti2026Power.stronglyMeasurable_expRegret}\leanok
\uses{def:bai_alg, cor:pb_gaussian_likelihood}
Let $\Xs$ be an action set and $\mathsf{Alg} = (\mathrm{alg},\rho)$ an identification algorithm
with budget $T\ge1$ on $\Xs$ (Definition~\ref{def:bai_alg}). Define, for $\theta\in\R^d$,
\[
  K(\theta) := \Prob_0^T.\mathrm{withDensity}\big(\mathrm{ENNReal.ofReal}\circ\ell_\theta\big),
  \qquad \kappa(\theta) := K(\theta)\otimes\rho .
\]
Then $K$ is a Markov kernel from $\R^d$ to histories of length $T$ with $K(\theta) = \Prob_\theta^T$
for every $\theta$, and $\kappa$ is a Markov kernel from $\R^d$ to pairs (history,
recommendation) with $\kappa(\theta) = \Prob_\theta^T\otimes\rho = \Prob_\theta$, the law of
$(\hist_T,\rec)$ of Definition~\ref{def:bai_alg}, for \emph{every} $\theta$. In particular
$\kappa$ satisfies the hypothesis of Lemma~\ref{lem:pac_expected_regret} for every prior
$\pi$, and for every bounded measurable $g$ of $(\theta,\hist_T,\rec)$ the map
$\theta\mapsto\E_\theta[g(\theta,\hist_T,\rec)] = \int g(\theta,\cdot)\,d\kappa(\theta)$ is
measurable.
\end{lemma}

\begin{proof}
\leanok
\uses{cor:pb_gaussian_likelihood}
$K$ is \texttt{Kernel.withDensity} of the constant kernel $\theta\mapsto\Prob_0^T$ with the
density $(\theta,h)\mapsto\mathrm{ENNReal.ofReal}(\ell_\theta(h))$, which is jointly measurable
by Corollary~\ref{cor:pb_gaussian_likelihood}; this joint measurability is exactly what
\texttt{Kernel.withDensity} requires to produce a kernel, and $K(\theta) = \Prob_\theta^T$ by
the same corollary (\texttt{Kernel.withDensity\_apply}). Each $K(\theta)$ is a probability
measure, so $K$ is a Markov kernel. Then $\kappa := K\otimes\rho'$ with
$\rho' := \rho.\mathrm{prodMkLeft}$ (the kernel $(\theta,h)\mapsto\rho(h)$) is a Markov kernel
(\texttt{Kernel.compProd} of Markov kernels) with
$\kappa(\theta) = K(\theta)\otimes\rho = \Prob_\theta^T\otimes\rho$
(\texttt{Kernel.compProd\_apply\_eq\_compProd\_sectR}, \texttt{Kernel.prodMkLeft\_apply}),
which is $\Prob_\theta$ by Definition~\ref{def:bai_alg} and the uniqueness of the law of the
history. The measurability of $\theta\mapsto\int g\,d\kappa(\theta)$ is the measurability of
integrals against a kernel (\texttt{Kernel.measurable\_lintegral},
\texttt{Measurable.lintegral\_kernel\_prod\_right}, and
\texttt{StronglyMeasurable.integral\_kernel\_prod\_right} for real-valued bounded $g$).
\end{proof}

\textbf{Lean remark.} This is the construction announced in the Lean remark of
Lemma~\ref{lem:pac_expected_regret}: the measurability of
$\theta\mapsto\mathrm{trajMeasure}(\mathrm{alg},\nu_\theta)$ is never needed, because the kernel
is built from the single reference law $\Prob_0^T$ and the explicit density $\ell_\theta$.

\section{The joint law of the parameter, the history and the recommendation}
\label{sec:pre_bayes_joint}

\begin{definition}[Bayesian joint law and posterior quantities]\label{def:bayes_joint}
\lean{Maiti2026Power.noisePairLaw, Maiti2026Power.pairKernel_eq_withDensity, Maiti2026Power.bayesJointLaw, Maiti2026Power.bayesJointLaw_eq_withDensity, ProbabilityTheory.gaussPrior, ProbabilityTheory.postPrec, ProbabilityTheory.postCov, ProbabilityTheory.postMean, ProbabilityTheory.posDef_postPrec, ProbabilityTheory.posDef_postCov, ProbabilityTheory.toEuclideanCLM_postPrec_postMean, ProbabilityTheory.sub_toEuclideanCLM_postMean}\leanok
\uses{def:bai_alg, cor:pb_gaussian_likelihood, lem:pb_kernel_of_likelihood, lem:loewner_inv, lem:eigen_quadratic_bounds}
Let $d\ge1$, $\Xs$ an action set, $\mathsf{Alg} = (\mathrm{alg},\rho)$ an identification algorithm
with budget $T\ge1$ on $\Xs$ (Definition~\ref{def:bai_alg}) and $\sigma>0$. Let
$\pi := \Normal(0,\sigma^2I_d)$ (the prior). The \emph{joint law} of $(\theta,\hist_T,\rec)$ is
the probability measure on $\R^d\times\mathcal{H}_T\times\Xs$
\[
  Q := \pi\otimes\kappa,\qquad \kappa(\theta) = \Prob_\theta^T\otimes\rho = \Prob_\theta ,
\]
where $\kappa$ is the Markov kernel of Lemma~\ref{lem:pb_kernel_of_likelihood} (built from
$\Prob_\theta^T = \Prob_0^T.\mathrm{withDensity}(\ell_\theta)$,
Corollary~\ref{cor:pb_gaussian_likelihood}). Equivalently, with $Q_0 := \Prob_0^T\otimes\rho$
the law of (history, recommendation) under pure noise,
$\kappa(\theta) = Q_0.\mathrm{withDensity}((h,x)\mapsto\ell_\theta(h))$ and
\[
  Q = (\pi\times Q_0).\mathrm{withDensity}\big((\theta,h,x)\mapsto\ell_\theta(h)\big) .
\]
We write $\E$ for the expectation under $Q$ and, for
a history $h$ of length $T$,
\[
  V_T(h) := \big(\sigma^{-2}I_d+S_T(h)\big)^{-1}\in\PD,\qquad m_T(h) := V_T(h)\,b_T(h)\in\R^d ,
\]
so that $(\sigma^{-2}I + S_T)m_T = b_T$ and $b_T - S_Tm_T = \sigma^{-2}m_T$.
\end{definition}

\textbf{Lean remark.} The pure-noise laws are \texttt{noiseHistLaw} and \texttt{noisePairLaw},
the kernel is \texttt{pairKernel}, the joint law is \texttt{bayesJointLaw A T σ} and its
product-with-density form is \texttt{bayesJointLaw\_eq\_withDensity}. The posterior quantities
are defined for a general tilt $(b, S)$: \texttt{postPrec σ S} $= \sigma^{-2}I + S$,
\texttt{postCov σ S} its inverse and \texttt{postMean σ b S} $= V b$; they are measurable
functions of the history through $b_T$, $S_T$, but this measurability is never needed (see the
Lean remark after Lemma~\ref{lem:posterior_second_moment}).

\begin{lemma}[Fubini for the joint law]\label{lem:pb_joint_fubini}
\lean{Maiti2026Power.integral_bayesJointLaw, Maiti2026Power.integrable_bayesJointLaw_iff, Maiti2026Power.integral_likelihood_noisePairLaw, Maiti2026Power.integrable_likelihood_noisePairLaw, Maiti2026Power.integrable_likelihood_mul_prod, Maiti2026Power.integral_likelihood_mul_prod, Maiti2026Power.integral_indicator_mul_simpleRegret_bayesJointLaw}\leanok
\uses{def:bayes_joint, cor:pb_gaussian_likelihood}
For every measurable $f : \R^d\times\mathcal{H}_T\times\Xs\to[0,\infty]$,
\[
  \E[f(\theta,\hist_T,\rec)]
  = \int\!\!\int\!\!\int f(\theta,h,x)\,\rho(h,dx)\,\Prob_\theta^T(dh)\,\pi(d\theta)
  = \int\!\!\int\ell_\theta(h)\,f(\theta,h,x)\,Q_0(dh,dx)\,\pi(d\theta)
  = \int\!\!\int\ell_\theta(h)\,f(\theta,h,x)\,\pi(d\theta)\,Q_0(dh,dx),
\]
and the same holds for real-valued $f$ that are $Q$-integrable, i.e.\ such that
$\ell\cdot f$ is $\pi\times Q_0$-integrable. In particular $\int\ell_\theta\,dQ_0 = 1$ for every
$\theta$, the marginal law of $\theta$ is $\pi$, the conditional law
of $(\hist_T,\rec)$ given $\theta$ is $\Prob_\theta$, and for a function $g$ of $\theta$ alone,
$\int\ell_\theta(h)g(\theta)\,d(\pi\times Q_0) = \int g\,d\pi$.
\end{lemma}

\begin{proof}
\leanok
\uses{cor:pb_gaussian_likelihood}
The first equality is Tonelli/Fubini for composition-products (Mathlib
\texttt{Measure.lintegral\_compProd}, \texttt{Measure.integral\_compProd}), the second is the
density of Corollary~\ref{cor:pb_gaussian_likelihood} and the product form of $Q$
(\texttt{integral\_withDensity\_eq\_integral\_toReal\_smul}), the third is Fubini for the product
$\pi\times Q_0$ (\texttt{integral\_prod}, \texttt{integral\_prod\_symm}). Since
$\kappa(\theta)$ is a probability measure, $\int\ell_\theta\,dQ_0 = \kappa(\theta)(\text{all}) = 1$.
\end{proof}

\section{Posterior mean and posterior second moment}
\label{sec:pre_bayes_posterior}

\begin{lemma}[Stein's identity for functions of exponential growth]\label{lem:gauss_ibp_exp}
\lean{ProbabilityTheory.integral_sub_mul_gaussianReal_of_integrable, ProbabilityTheory.integral_mul_gaussianReal_zero_one_of_integrable, ProbabilityTheory.integrable_exp_mul_abs_gaussianReal, ProbabilityTheory.IsGaussian.integrable_of_abs_le_mul_exp, ProbabilityTheory.IsGaussian.integrable_norm_mul_of_abs_le_mul_exp, ProbabilityTheory.IsGaussian.integrable_inner_mul_of_abs_le_mul_exp, ProbabilityTheory.IsGaussian.integrable_fderiv_apply_of_norm_fderiv_le_mul_exp, ProbabilityTheory.integral_apply_mul_stdGaussian_of_le_exp, ProbabilityTheory.integral_inner_mul_stdGaussian_euclidean_of_le_exp, ProbabilityTheory.integral_inner_mul_stdGaussian_of_le_exp, ProbabilityTheory.integral_inner_mul_multivariateGaussian_of_le_exp, ProbabilityTheory.norm_toLp_insertNth_le, ProbabilityTheory.abs_apply_toLp_insertNth_le}\leanok
\uses{lem:gauss_ibp, lem:psf_ibp_1d, def:pg_gaussian_vector}
Let $F : \R^d\to\R$ be $C^1$ with $\abs{F(x)} \le Ae^{B\norm x}$ and
$\norm{DF(x)} \le Ae^{B\norm x}$ for all $x$ (some constants $A, B$). Then for
$g\sim\Normal(0,I_d)$ and every $a\in\R^d$, $\E[\ip{a}{g}F(g)] = \E[DF(g)\,a]$, and for
$X\sim\Normal(0,S)$ with $S\in\PSD$, $\E[\ip{a}{X}F(X)] = \E[DF(X)\,(Sa)]$.
\end{lemma}

\begin{proof}
\leanok
\uses{lem:gauss_ibp, lem:psf_ibp_1d}
The one-dimensional identity $\E[\xi f(\xi)] = \E[f'(\xi)]$ (Lemma~\ref{lem:psf_ibp_1d}) only
needs the integrability of $f$, $\xi f(\xi)$ and $f'$, which holds for functions of exponential
growth since $e^{B\abs t}$ has all Gaussian moments. The multivariate identity follows by the
coordinate-wise Fubini argument of Lemma~\ref{lem:gauss_ibp}: the section
$t\mapsto F(x + te_i)$ has growth $Ae^{B\norm{x}}e^{B\abs t}$, and the integrability of
$\ip{a}{g}F(g)$ and $DF(g)a$ under a Gaussian measure is Fernique's theorem.
\end{proof}

\begin{lemma}[Gaussian tilt: derivative and growth]\label{lem:pb_complete_square_identity}
\lean{ProbabilityTheory.gaussTilt, ProbabilityTheory.gaussTilt_pos, ProbabilityTheory.gaussTilt_le, ProbabilityTheory.hasFDerivAt_gaussTilt, ProbabilityTheory.contDiff_gaussTilt, ProbabilityTheory.hasFDerivAt_affine_mul_gaussTilt, ProbabilityTheory.abs_affine_mul_gaussTilt_le, ProbabilityTheory.norm_fderiv_affine_mul_gaussTilt_le, ProbabilityTheory.inner_toEuclideanCLM_symm, ProbabilityTheory.inner_toEuclideanCLM_nonneg}\leanok
Let $S\in\PSD$, $b\in\R^d$ and $\ell(\theta) := \exp(\ip{b}{\theta}-\tfrac12\theta^\top S\theta)$.
Then $\ell > 0$, $\ell(\theta) \le e^{\norm b\norm\theta}$, $\ell$ is $C^1$ with
$D\ell(\theta) = \ell(\theta)\,\ip{b - S\theta}{\cdot}$, and for $c\in\R^d$, $r\in\R$ the
function $\theta\mapsto(\ip{c}{\theta}+r)\ell(\theta)$ and its derivative have exponential
growth: they are bounded by $A e^{(\norm b+2)\norm\theta}$ with
$A = (\norm c+\abs r)(\norm b+\norm S+1)+\norm c$.
\end{lemma}

\begin{proof}
\leanok
$\theta^\top S\theta\ge0$ gives the bound; the derivative is the chain rule, and the growth
bounds use $\norm\theta \le e^{\norm\theta}$.
\end{proof}

\begin{lemma}[Stein's identity for the affine tilt]\label{lem:pb_gaussian_integrals}
\lean{ProbabilityTheory.integral_inner_mul_affine_mul_gaussTilt, ProbabilityTheory.integrable_affine_mul_affine_mul_gaussTilt, ProbabilityTheory.integrable_affine_mul_gaussTilt, ProbabilityTheory.integrable_gaussTilt, ProbabilityTheory.integrable_gaussTilt_smul, ProbabilityTheory.integral_inner_mul_gaussTilt, ProbabilityTheory.integral_gaussTilt_pos}\leanok
\uses{lem:gauss_ibp_exp, lem:pb_complete_square_identity}
Let $\sigma>0$, $\pi = \Normal(0,\sigma^2I_d)$, $S\in\PSD$, $b\in\R^d$ and $\ell$ as above.
For all $a, c\in\R^d$ and $r\in\R$,
\[
  \int\ip{a}{\theta}\,(\ip{c}{\theta}+r)\,\ell(\theta)\,\pi(d\theta)
  = \sigma^2\int\big(\ip{c}{a} + (\ip{c}{\theta}+r)\ip{b-S\theta}{a}\big)\ell(\theta)\,\pi(d\theta).
\]
All the functions $\theta\mapsto(\ip{a}{\theta}+r)(\ip{c}{\theta}+s)\ell(\theta)$ are
$\pi$-integrable, $Z := \int\ell\,d\pi\in(0,\infty)$, and
$\int\ip{a}{\theta}\ell(\theta)\,\pi(d\theta) = \ip{a}{\int\ell(\theta)\theta\,\pi(d\theta)}$.
\end{lemma}

\begin{proof}
\leanok
\uses{lem:gauss_ibp_exp, lem:pb_complete_square_identity}
Lemma~\ref{lem:gauss_ibp_exp} for $\Normal(0,\sigma^2I)$ applied to
$F = (\ip{c}{\cdot}+r)\ell$, whose derivative is computed in
Lemma~\ref{lem:pb_complete_square_identity}; $Sa = \sigma^2a$ here.
\end{proof}

\begin{lemma}[Gaussian posterior identities]\label{lem:gaussian_complete_square}
\lean{ProbabilityTheory.integral_gaussTilt_smul_eq, ProbabilityTheory.integral_inner_sub_postMean_mul_gaussTilt, ProbabilityTheory.integral_inner_sub_mul_inner_sub_mul_gaussTilt_add, ProbabilityTheory.integral_inner_single_sub_mul_inner_sub_mul_gaussTilt, ProbabilityTheory.integral_norm_sq_sub_postMean_mul_gaussTilt, ProbabilityTheory.integral_norm_sq_mul_gaussTilt, ProbabilityTheory.integrable_inner_sub_mul_inner_sub_mul_gaussTilt, ProbabilityTheory.integrable_inner_sub_mul_gaussTilt}\leanok
\uses{lem:pb_gaussian_integrals, def:bayes_joint}
Let $\sigma>0$, $S\in\PSD$, $b\in\R^d$, $V := (\sigma^{-2}I_d+S)^{-1}$, $m := Vb$,
$\pi := \Normal(0,\sigma^2I_d)$, $\ell(\theta) := \exp(\ip{b}{\theta}-\tfrac12\theta^\top S\theta)$
and $Z := \int\ell\,d\pi > 0$. Then
\[
  \int\ell(\theta)\,\theta\,\pi(d\theta) = Z\,m,\qquad
  \int\ip{c}{\theta-m}\,\ell(\theta)\,\pi(d\theta) = 0\ \ (c\in\R^d),
\]
\[
  \int\norm{\theta-m}^2\ell(\theta)\,\pi(d\theta) = \tr(V)\,Z,\qquad
  \int\norm{\theta}^2\ell(\theta)\,\pi(d\theta) = (\tr V + \norm m^2)\,Z .
\]
(Equivalently, $\ell\,d\pi/Z$ is the Gaussian $\Normal(m,V)$; we only need these moments.)
\end{lemma}

\begin{proof}
\leanok
\uses{lem:pb_gaussian_integrals, def:bayes_joint}
\emph{Mean.} Lemma~\ref{lem:pb_gaussian_integrals} with $c = 0$, $r = 1$: writing
$u := \int\ell(\theta)\theta\,d\pi$, for every $a$,
$\ip{a}{u} = \sigma^2(\ip{b}{a}Z - \ip{a}{Su})$ (using $\ip{S\theta}{a} = \ip{\theta}{Sa}$).
Hence $u + \sigma^2Su = \sigma^2Zb$, i.e.\ $(\sigma^{-2}I+S)u = Zb$ and $u = Zm$. The second
identity is $\ip{c}{u} - \ip{c}{m}Z = 0$.

\emph{Second moment.} Let $\Phi(a,c) := \int\ip{a}{\theta-m}\ip{c}{\theta-m}\ell\,d\pi$.
Lemma~\ref{lem:pb_gaussian_integrals} with $r = -\ip{c}{m}$ (so that
$\ip{c}{\theta}+r = \ip{c}{\theta-m}$), together with
$b - S\theta = \sigma^{-2}m - S(\theta-m)$ (Definition~\ref{def:bayes_joint}) and the mean
identity, gives $\Phi(a,c) + \sigma^2\Phi(Sa,c) = \sigma^2\ip{c}{a}Z$ for all $a, c$. With
$a := \sigma^{-2}Ve_i$ one has $a + \sigma^2Sa = \sigma^2(\sigma^{-2}I+S)a = e_i$, so
$\Phi(e_i,c) = \ip{c}{Ve_i}Z$; taking $c = e_i$ and summing over $i$ gives
$\int\norm{\theta-m}^2\ell\,d\pi = \tr(V)Z$. Finally
$\norm\theta^2 = \norm{\theta-m}^2 + 2\ip{m}{\theta-m} + \norm m^2$ and the cross term
integrates to $0$.
\end{proof}

\textbf{Lean remark.} \texttt{PosteriorGaussian.lean} contains this section for the general tilt
$(b, S)$ (\texttt{gaussTilt b S}, \texttt{postCov σ S}, \texttt{postMean σ b S}); it is applied
to $b = b_T(h)$, $S = S_T(h)$ for each fixed history $h$ through
\texttt{Maiti2026Power.likelihood\_eq\_gaussTilt}.

\begin{lemma}[Posterior mean identity]\label{lem:posterior_mean}
\lean{Maiti2026Power.postU, Maiti2026Power.stronglyMeasurable_postU, Maiti2026Power.integrable_likelihood_mul_inner_prod, Maiti2026Power.integrable_norm_postU, Maiti2026Power.integral_inner_bayesJointLaw_le}\leanok
\uses{def:bayes_joint, lem:pb_joint_fubini, lem:gaussian_complete_square}
Let $\Xs \subseteq \ball_d$ and define, for a history $h$,
$u_T(h) := \int\ell_\theta(h)\,\theta\,\pi(d\theta) = Z_T(h)\,m_T(h)$ with
$Z_T(h) := \int\ell_\theta(h)\,\pi(d\theta)$; $u_T$ is a measurable function of $h$ with
$\int\norm{u_T}\,dQ_0 < \infty$. Then
\[
  \E\ip{\rec}{\theta} = \int\ip{x}{u_T(h)}\,Q_0(dh,dx) \le \int\norm{u_T(h)}\,Q_0(dh,dx) .
\]
\end{lemma}

\begin{proof}
\leanok
\uses{lem:pb_joint_fubini, lem:gaussian_complete_square}
By Lemma~\ref{lem:pb_joint_fubini} with the $\theta$-integral innermost (the integrand
$\ell_\theta(h)\ip{x}{\theta}$ is $\pi\times Q_0$-integrable, being dominated by
$\ell_\theta(h)\norm\theta$ whose integral is $\E_\pi\norm\theta$),
$\E\ip{\rec}{\theta} = \int\int\ell_\theta(h)\ip{x}{\theta}\,\pi(d\theta)\,Q_0(dh,dx)
= \int\ip{x}{u_T(h)}\,Q_0(dh,dx)$ (Lemma~\ref{lem:pb_gaussian_integrals}), and
$\ip{x}{u_T(h)} \le \norm{u_T(h)}$ since $\norm x\le1$. Measurability of $u_T$ is that of a
parametric integral (Mathlib \texttt{StronglyMeasurable.integral\_prod\_right'}).
\end{proof}

\begin{lemma}[Posterior second moment]\label{lem:posterior_second_moment}
\lean{Maiti2026Power.postZ, Maiti2026Power.postN, Maiti2026Power.stronglyMeasurable_postZ, Maiti2026Power.stronglyMeasurable_postN, Maiti2026Power.postZ_pos, Maiti2026Power.norm_postU_sq_le, Maiti2026Power.norm_postU_le, Maiti2026Power.integrable_postZ, Maiti2026Power.integral_postZ, Maiti2026Power.integrable_postN, Maiti2026Power.integral_postN, Maiti2026Power.integral_norm_postU_le}\leanok
\uses{def:bayes_joint, lem:pb_joint_fubini, lem:gaussian_complete_square, lem:trace_inv_lower}
Let $\Xs \subseteq \ball_d$ and $N_T(h) := \int\norm\theta^2\ell_\theta(h)\,\pi(d\theta)$. For
every history $h$, with $c := d^2/(d\sigma^{-2}+T)$,
\[
  \norm{u_T(h)}^2 = Z_T(h)^2\norm{m_T(h)}^2 \le Z_T(h)\big(N_T(h) - c\,Z_T(h)\big),
\]
and $\int Z_T\,dQ_0 = 1$, $\int N_T\,dQ_0 = \E_\pi\norm\theta^2 = \sigma^2d$. Consequently
\[
  \int\norm{u_T(h)}\,Q_0(dh,dx)\ \le\ \sqrt{\sigma^2d - \frac{d^2}{d\sigma^{-2}+T}} .
\]
\end{lemma}

\begin{proof}
\leanok
\uses{lem:pb_joint_fubini, lem:gaussian_complete_square, lem:trace_inv_lower}
By Lemma~\ref{lem:gaussian_complete_square} for the fixed history $h$,
$u_T = Z_Tm_T$ and $N_T = (\tr V_T + \norm{m_T}^2)Z_T$, so
$Z_T\norm{m_T}^2 = N_T - \tr(V_T)Z_T \le N_T - cZ_T$ by Lemma~\ref{lem:trace_inv_lower}; multiply
by $Z_T$. The integrals of $Z_T$ and $N_T$ are Lemma~\ref{lem:pb_joint_fubini} (Fubini, the
$\theta$-integral outermost). Finally, by the Cauchy--Schwarz inequality in $L^2(Q_0)$ applied
to $\sqrt{Z_T}$ and $\sqrt{N_T - cZ_T}$,
$\int\norm{u_T}\,dQ_0 \le \int\sqrt{Z_T}\sqrt{N_T-cZ_T}\,dQ_0
\le \big(\int Z_T\,dQ_0\big)^{1/2}\big(\int(N_T - cZ_T)\,dQ_0\big)^{1/2} = \sqrt{\sigma^2d - c}$.
\end{proof}

\textbf{Lean remark.} This route never integrates $m_T$ or $\tr V_T$ over histories, only
$Z_T$, $N_T$ and $u_T$, which are parametric integrals in $h$ (measurable by
\texttt{StronglyMeasurable.integral\_prod\_right'}); the outline's integrability lemma
$\E\norm{b_T}^2 < \infty$ and the measurability of $h\mapsto V_T(h)$ (matrix inversion) are
not needed.

\section{The unit-ball lower bound}
\label{sec:pre_bayes_ball}

\begin{lemma}[Lower bound on the posterior trace]\label{lem:trace_inv_lower}
\lean{Maiti2026Power.le_trace_postCov_histS, Maiti2026Power.trace_histS_le, Maiti2026Power.posSemidef_histS}\leanok
\uses{def:bayes_joint, lem:trace_inv_amhm}
If $\Xs\subseteq\ball_d$, then for every history $h$, $\tr V_T(h)\ge\frac{d^2}{d\sigma^{-2}+T}$.
\end{lemma}

\begin{proof}
\leanok
\uses{lem:trace_inv_amhm}
$\tr S_T(h) = \sum_{t<T}\norm{x_t}^2\le T$, so $\tr(\sigma^{-2}I+S_T(h))\le d\sigma^{-2}+T$, and
Lemma~\ref{lem:trace_inv_amhm} applied to the positive definite matrix $\sigma^{-2}I+S_T(h)$
gives $\tr V_T(h)\ge d^2/\tr(\sigma^{-2}I+S_T(h))\ge d^2/(d\sigma^{-2}+T)$.
\end{proof}

\begin{lemma}[Expected norm of a Gaussian vector, lower bound]\label{lem:pb_norm_lower}
\lean{Maiti2026Power.le_integral_norm_gaussPrior, Maiti2026Power.integral_norm_sq_gaussPrior, Maiti2026Power.integrable_norm_gaussPrior, Maiti2026Power.integrable_norm_sq_gaussPrior}\leanok
\uses{lem:gauss_norm_lower, lem:gauss_norm_moments}
If $\theta\sim\Normal(0,\sigma^2I_d)$ then $\E\norm\theta^2 = \sigma^2d$ and
$\E\norm\theta\ge\sqrt{2/\pi}\,\sigma\sqrt d$.
\end{lemma}

\begin{proof}
\leanok
\uses{lem:gauss_norm_lower, lem:gauss_norm_moments}
Lemma~\ref{lem:gauss_norm_moments}(i) and Lemma~\ref{lem:gauss_norm_lower} with the weights
$c_i = 1/\sqrt d$ and $S = \sigma^2I_d$.
\end{proof}

\begin{lemma}[Bayesian simple regret on the ball]\label{lem:bayes_regret_ball}
\lean{Maiti2026Power.integral_simpleRegret_bayesJointLaw_ge}\leanok
\uses{def:bayes_joint, def:regret, lem:posterior_mean, lem:posterior_second_moment, lem:pb_norm_lower}
Let $\Xs = \ball_d$. Under the joint law $Q$ of Definition~\ref{def:bayes_joint},
\[
  \E[r(\rec,\theta)] = \E\norm\theta-\E\ip{\rec}{\theta}
  \ \ge\ \sqrt{2/\pi}\,\sigma\sqrt d-\sqrt{\sigma^2d-\frac{d^2}{d\sigma^{-2}+T}} .
\]
\end{lemma}

\begin{proof}
\leanok
\uses{lem:posterior_mean, lem:posterior_second_moment, lem:pb_norm_lower}
On the unit ball, $\max_{x\in\ball_d}\ip{x}{\theta} = \norm\theta$, so
$r(\rec,\theta) = \norm\theta-\ip{\rec}{\theta}$ (Lemma~\ref{lem:block_ball_basic}). The marginal
of $\theta$ is $\pi$, so $\E\norm\theta \ge \sqrt{2/\pi}\sigma\sqrt d$
(Lemma~\ref{lem:pb_norm_lower}), and
$\E\ip{\rec}{\theta} \le \int\norm{u_T}\,dQ_0 \le \sqrt{\sigma^2d - d^2/(d\sigma^{-2}+T)}$
by Lemmas~\ref{lem:posterior_mean} and~\ref{lem:posterior_second_moment}.
\end{proof}

\begin{theorem}[Simple-regret lower bound on the unit ball]\label{thm:ball_simple_regret_lower}
\lean{Maiti2026Power.simpleRegret_unitBall_nonneg, Maiti2026Power.simpleRegret_unitBall_le, Maiti2026Power.expRegret_nonneg, Maiti2026Power.expRegret_le, Maiti2026Power.integrable_expRegret, Maiti2026Power.integrable_simpleRegret_bayesJointLaw, Maiti2026Power.integral_norm_sq_bayesJointLaw, Maiti2026Power.integral_indicator_mul_simpleRegret_bayesJointLaw_ge, Maiti2026Power.le_measureReal_norm_le_gaussPrior, Maiti2026Power.exists_norm_le_and_le_expRegret, Maiti2026Power.exists_norm_le_and_le_expRegret_unitBall}\leanok
\uses{def:bayes_joint, def:regret, lem:bayes_regret_ball, lem:pb_joint_fubini, lem:gauss_norm_moments}
Let $d\ge1$, $T\ge1$ and $\mathsf{Alg}$ be any identification algorithm with budget $T$ on
$\Xs = \ball_d$. With the prior $\pi = \Normal(0,\sigma^2I_d)$, $\sigma^2 := d/(4T)$:
\begin{enumerate}
  \item[(i)] $\E[r(\rec,\theta)]\ge\dfrac{d}{8\sqrt T}$ under the joint law $Q$;
  \item[(ii)] there exists $\theta\in\R^d$ with $\norm\theta\le\dfrac{10\,d}{\sqrt T}$ and
        $\E_\theta[r(\rec,\theta)]\ge\dfrac{3d}{40\sqrt T}$.
\end{enumerate}
\end{theorem}

\begin{proof}
\leanok
\uses{lem:bayes_regret_ball, lem:pb_joint_fubini, lem:gauss_norm_moments, lem:pb_kernel_of_likelihood}
(i) With $\sigma^2 = d/(4T)$ we have $d\sigma^{-2} = 4T$, hence
$\frac{d^2}{d\sigma^{-2}+T} = \frac{d^2}{5T}$ and $\sigma^2d-\frac{d^2}{5T} = \frac{d^2}{4T}-\frac{d^2}{5T} = \frac{d^2}{20T}$,
while $\sqrt{2/\pi}\,\sigma\sqrt d = \sqrt{2/\pi}\,\frac{d}{2\sqrt T} \ge 0.79\cdot\frac{d}{2\sqrt T}$
(as $\pi < 3.15$). Lemma~\ref{lem:bayes_regret_ball} gives
\[
  \E[r(\rec,\theta)]\ge\frac{d}{\sqrt T}\Big(\frac{0.79}{2}-\frac{1}{4.47}\Big)
  \ge\frac{d}{\sqrt T}\,(0.395-0.2238) \ge \frac{d}{8\sqrt T} ,
\]
using $4.47^2 \le 20$.

(ii) Let $R := \frac{10d}{\sqrt T} = 20\sigma\sqrt d$.
Since $0 \le r(\rec,\theta)\le2\norm\theta$ on the ball and
$\norm\theta\indic\{\norm\theta>R\}\le\norm\theta^2/R$,
\[
  \E\big[r(\rec,\theta)\indic\{\norm\theta>R\}\big]\le\frac{2\,\E\norm\theta^2}{R}
  = \frac{2\sigma^2d}{R} = \frac{d}{20\sqrt T} ,
\]
so $\E[r\indic\{\norm\theta\le R\}] \ge \frac{d}{8\sqrt T}-\frac{d}{20\sqrt T} = \frac{3d}{40\sqrt T}$.
By Lemma~\ref{lem:pb_joint_fubini}, $\E[r\indic\{\norm\theta\le R\}] = \int_{\norm\theta\le R}\E_\theta[r(\rec,\theta)]\,\pi(d\theta)$,
where $\theta\mapsto\E_\theta[r(\rec,\theta)]$ is measurable
(Lemma~\ref{lem:pb_kernel_of_likelihood}, $r$ being bounded by $2\norm\theta$).
If every $\theta$ with $\norm\theta\le R$ had $\E_\theta[r]<\frac{3d}{40\sqrt T}$, the
nonnegative function $g := \indic\{\norm\theta\le R\}\,(\frac{3d}{40\sqrt T} - \E_\theta[r])$
would be positive on $\{\norm\theta\le R\}$, a set of positive $\pi$-measure (by Markov's
inequality $\pi(\norm\theta>R)\le\sigma^2d/R^2 = 1/400$), so $\int g\,d\pi > 0$; but
$\int g\,d\pi = \frac{3d}{40\sqrt T}\pi(\norm\theta\le R) - \E[r\indic\{\norm\theta\le R\}] \le 0$,
a contradiction. Hence some $\theta$ with $\norm\theta\le R$ satisfies
$\E_\theta[r(\rec,\theta)]\ge\frac{3d}{40\sqrt T}$.
\end{proof}

\begin{remark}[Constants]\label{rem:pb_constants}
The outline claimed the constant $0.13$ in (i), from
$\frac{1}{2\sqrt2}-\frac{1}{\sqrt{20}} = 0.12995\ldots$, which is slightly \emph{below} $0.13$; we
use $1/8$ instead (the sharper Lemma~\ref{lem:pb_norm_lower} would even give $0.17$). The
truncation radius was changed from $16\sigma\sqrt d$ to $20\sigma\sqrt d$
so that Corollary~\ref{cor:ball_pac_lower} keeps the outline's headline constants
($\delta\le1/500$, $T\ge d^2/(1000\eps^2)$): with $R = k\sigma\sqrt d$ the final coefficient is
$\frac18-\frac1k-k\delta$, which for $\delta = 1/500$ is maximized near $k = 22$; $k = 20$ gives
$7/200$. The choice $\sigma^2 = d/(4T)$ is not optimal either but keeps the arithmetic simple.
\end{remark}

\begin{corollary}[PAC lower bound on the unit ball]\label{cor:ball_pac_lower}
\lean{Maiti2026Power.unitBall_le_budget_of_isPAC, Maiti2026Power.pairKernel_zero_eq, Maiti2026Power.not_isPAC_unitBall_of_isFixedBudget_zero}\leanok
\uses{def:pac, lem:pac_expected_regret, lem:pb_kernel_of_likelihood, thm:ball_simple_regret_lower}
Let $d\ge2$, $\eps>0$ and $\delta\le1/500$. Every $(\eps,\delta)$-PAC identification algorithm
on $\ball_d$ with budget $T$ satisfies
\[
  T\ge\frac{49\,d^2}{40000\,\eps^2}\ge\frac{d^2}{1000\,\eps^2} .
\]
\end{corollary}

\begin{proof}
\leanok
\uses{lem:pac_expected_regret, lem:pb_kernel_of_likelihood, thm:ball_simple_regret_lower}
If $T = 0$ the law of the recommendation does not depend on $\theta$, and the two instances
$\pm2\eps e_1$ have regrets summing to $4\eps$ at every arm, so the algorithm cannot be
$(\eps,\delta)$-PAC for $\delta<1/2$ (as in Lemma~\ref{lem:singular_design_fails}).
Let $T\ge1$ and let $\theta$ be given by Theorem~\ref{thm:ball_simple_regret_lower}(ii). On
the ball the regret lies in $[0, 2\norm\theta]$, so Lemma~\ref{lem:pac_expected_regret}
applied to the law $\Prob_\theta = \kappa(\theta)$ of Lemma~\ref{lem:pb_kernel_of_likelihood}
gives $\E_\theta[r(\rec,\theta)]\le\eps+2\delta\norm\theta\le\eps+\frac{20\delta d}{\sqrt T}$
(if $2\norm\theta < \eps$ this is trivial). Combining
with $\E_\theta[r]\ge\frac{3d}{40\sqrt T}$,
\[
  \Big(\frac{3}{40}-20\delta\Big)\frac{d}{\sqrt T}\le\eps,\qquad
  \frac{3}{40}-20\delta\ge\frac{3}{40}-\frac{1}{25} = \frac{7}{200},
\]
hence $\sqrt T\ge\frac{7d}{200\,\eps}$ and $T\ge\frac{49d^2}{40000\eps^2}$; finally
$40000/49<1000$.
\end{proof}

\textbf{Lean remark.} The unit ball is \texttt{unitBall ι = Metric.closedBall 0 1} in
\texttt{EuclideanSpace ℝ ι}; it is compact and spans $\R^d$, hence an action set. The
expectation $\E_\theta[r(\rec,\theta)]$ is \texttt{expRegret T A θ}, the integral of a bounded
function under $\Prob_\theta^T\otimes\rho = \kappa(\theta)$ (\texttt{pairKernel A T θ}), and its
measurability in $\theta$ is Mathlib's \texttt{StronglyMeasurable.integral\_kernel\_prod\_right'}.
